% ============================================================
%  sec-2-npu-foundations.tex  —  2.2 NPU Computing Foundations
% ============================================================

\section{NPU Computing Foundations}
\label{sec:npu-foundations}


% ── 2.2.1 ────────────────────────────────────────────────────
\subsection{NPU Architecture Principles}
\label{subsec:npu-arch}


The most prevalent NPU computational substrate is the \textit{systolic array}:
a 2D grid of Processing Elements (PEs) arranged so that data flows
rhythmically through the array like a heartbeat, with each element
performing a MAC operation as data passes through. In a typical
$N \times N$ configuration:

\begin{itemize}
  \item Input activations stream west-to-east across rows.
  \item Weight matrices are pre-loaded north-to-south down columns.
  \item Partial sums accumulate southward and are written out at the bottom edge.
\end{itemize}

The key advantage is that each data element is read once from memory and
reused across many computations as it flows through the grid, dramatically
reducing memory bandwidth requirements relative to a GPU's scratchpad model.

NPU memory design is governed by the principle that physical distance from
compute correlates directly with energy cost, with the four-level hierarchy
summarised in \tabref{npu-memory}.

\begin{table}[htbp]
  \centering
  \caption{NPU on-chip memory hierarchy and access costs}
  \label{tab:npu-memory}
  \begin{tabularx}{\textwidth}{lllXX}
    \toprule
    \tableheader{Level & Type & Size & Access Cost & Purpose} \\
    \midrule
    PE Register File & Register & Bytes  & 1$\times$ (baseline) & Immediate MAC operands \\
    \rowcolor{TableRowOdd}
    PE Local Scratchpad & SRAM   & KB     & $\sim$2$\times$      & Local PE data \\
    Global Buffer    & SRAM   & 16–24 MB & $\sim$20$\times$     & Active weights and activations \\
    \rowcolor{TableRowOdd}
    Off-chip Memory  & HBM/DRAM & GB    & 200–1,000$\times$   & Full model weights \\
    \bottomrule
  \end{tabularx}
\end{table}


% ── 2.2.3 ────────────────────────────────────────────────────
\subsection{Specialised Execution Hardware}
\label{subsec:npu-hw}

A typical NPU execution pipeline consists of three specialised units
working in concert. The \textbf{Matrix Engine} (or MAC Array, MAA) is a
systolic array implementing the core GEMM operation, often fused with a
bias accumulator. The \textbf{Vector Engine} handles element-wise
operations: ReLU, sigmoid, tanh, layer normalisation. The \textbf{Dispatching
Unit} (DSU) routes non-zero input feature maps to the MAA, enabling
zero-skipping to accelerate sparse computations.

NPUs exploit \textit{operator fusion} to eliminate intermediate memory
round-trips. A Conv+Bias+ReLU sequence, for instance, executes as a single
hardware pass — saving two global memory write-reads compared to three
separate kernel launches as required on a GPU. For transformer models,
the entire multi-head attention block (Q$\times$K, Softmax, $\times$V) can
be fused into a single dispatched operation.


% ── 2.2.5 ────────────────────────────────────────────────────
\subsection{NPU Performance Metrics}
\label{subsec:npu-metrics}

\acrfull{tops} is the primary throughput metric for NPU compute capacity.
The Huawei Ascend 910 achieves approximately 256 TFLOPS (FP16) and 512 TOPS
in INT8 mode. \tabref{perf-metrics} provides a quantitative comparison of
inference performance across processor types.

\begin{table}[htbp]
  \centering
  \caption{Inference performance comparison: CPU, GPU, NPU}
  \label{tab:perf-metrics}
  \begin{tabularx}{\textwidth}{lXXX}
    \toprule
    \tableheader{Metric & CPU & GPU & NPU} \\
    \midrule
    Peak TOPS (INT8) & $<$0.001 & 30–624 & 256–512 \\
    \rowcolor{TableRowOdd}
    FP16 TFLOPS     & $<$0.1   & 100–989 & 256 \\
    Latency (ResNet-50) & 50–100 ms & 5–10 ms & 1–3 ms \\
    \rowcolor{TableRowOdd}
    Throughput (img/s) & 10–20 & 100–300 & 300–600 \\
    Energy efficiency & 1–5 GFLOPS/W & 10–30 GFLOPS/W & 38–60 TOPS/W \\
    \rowcolor{TableRowOdd}
    Memory bandwidth & 50–100 GB/s & 200–3,350 GB/s & 1,228 GB/s (Ascend 910) \\
    \bottomrule
  \end{tabularx}
\end{table}

Experimental benchmarks (KAIST) show that an NPU delivers 60\% faster
inference than a comparable GPU while consuming 44\% less power for
medium-complexity models. At high utilisation (MobileNetV2: 99.9\%), NPUs
achieve 143$\times$ better energy efficiency than CPU equivalents, with
only 0.15–0.7\% accuracy loss under INT8 quantisation.
